\section{Q\&A 精华：深度问答中的关键洞见}

\subsection{结构方法的偏差来源}

\textbf{Q}：为什么纯结构方法的生成分布存在严重偏差？

\textbf{A}：根本原因在于\textbf{结构数据的采样偏差}。实验解析结构的蛋白质（PDB 数据库中约 30 万条）倾向于以下特征：
\begin{itemize}
    \item \textbf{球状蛋白} (globular proteins)：紧凑、可溶于水溶液，便于实验操作
    \item \textbf{$\alpha$-helix 丰富}：这类结构元素在已解析蛋白质中被过度代表
\end{itemize}

EvoDiff 的论文中提供了详细的数据支持：EvoDiff 的生成蛋白质能更好地涵盖 $\alpha$-helix 之外的其他结构元素（如 $\beta$-sheet 等），而纯结构方法则倾向于过采样富含 $\alpha$-helix 的蛋白质。

\subsection{序列与结构融合的前景}

\textbf{Q}：EvoDiff 未来是否可以整合结构信息？

\textbf{A}：是的，这是蛋白质 ML 领域的活跃研究方向。具体的融合策略包括：
\begin{itemize}
    \item 将序列和结构映射到\textbf{公共表征空间}
    \item 构建\textbf{联合的序列-结构生成模型}
    \item 利用结构信息来\textbf{约束和引导}序列生成
\end{itemize}

两种数据模态是\textbf{互补的}：结构数据量小但提供精确的几何约束，序列数据量大且覆盖生物多样性的广度。将它们融合有望兼得两者的优势。

\subsection{本章小结}

Q\&A 环节揭示了两个重要的技术洞见：(1) 结构数据的系统性偏差是纯结构方法的根本局限；(2) 序列与结构的融合是实现更强大蛋白质设计能力的必经之路。


